% ============================================================
%  Capítulo 3 — Intervalos de Confianza
% ============================================================
\section{Intervalos de Confianza}

\subsection{Definición}

\begin{defi}[Intervalo de confianza]
  Un \textbf{intervalo de confianza} al nivel $(1-\alpha)$ para $\theta$ es
  un intervalo aleatorio $[L(\mathbf{X}),\, U(\mathbf{X})]$ tal que
  \[
    P_\theta\bigl(L(\mathbf{X}) \leq \theta \leq U(\mathbf{X})\bigr)
    \geq 1 - \alpha \quad \forall\, \theta \in \Theta.
  \]
\end{defi}

\begin{cuidado}
  El intervalo es \emph{aleatorio}, no el parámetro. Una vez observada la
  muestra, el intervalo calculado $[l, u]$ o contiene a $\theta$ o no lo
  contiene; no tiene sentido hablar de probabilidad en ese caso.
\end{cuidado}

\subsection{IC para la media con varianza conocida}

Sea $X_1,\ldots,X_n \iid \N(\mu, \sigma^2)$ con $\sigma^2$ conocida. Entonces
\[
  Z = \frac{\bar{X} - \mu}{\sigma/\sqrt{n}} \sim \N(0,1),
\]
y el IC al $(1-\alpha)$ para $\mu$ es
\[
  \bar{X} \pm z_{\alpha/2}\,\frac{\sigma}{\sqrt{n}},
\]
donde $z_{\alpha/2}$ es el cuantil $1-\alpha/2$ de la $\N(0,1)$.

\subsection{IC para la media con varianza desconocida}

Cuando $\sigma^2$ es desconocida, se reemplaza por $S^2$ y se usa la
distribución $t$ de Student:
\[
  T = \frac{\bar{X} - \mu}{S/\sqrt{n}} \sim t_{n-1}.
\]
El IC al $(1-\alpha)$ para $\mu$ es
\[
  \bar{X} \pm t_{n-1,\,\alpha/2}\,\frac{S}{\sqrt{n}}.
\]

\subsection{IC para la varianza}

\[
  \frac{(n-1)S^2}{\sigma^2} \sim \chi^2_{n-1}
  \implies
  IC_{1-\alpha}(\sigma^2) =
  \left[
    \frac{(n-1)S^2}{\chi^2_{n-1,\,\alpha/2}},\;
    \frac{(n-1)S^2}{\chi^2_{n-1,\,1-\alpha/2}}
  \right].
\]
